% Generated draft: requires mathematical and source review.
\documentclass[11pt]{article}
\usepackage[margin=25mm]{geometry}
\usepackage{fontspec}
\setmainfont{Latin Modern Roman}
\usepackage{amsmath,amssymb,mathtools,unicode-math}
\setmathfont{Latin Modern Math}
\usepackage{enumitem,needspace,xurl,hyperref}
\setlength{\parindent}{0pt}
\setlength{\parskip}{5pt}
\setlength{\emergencystretch}{3em}
\newcommand{\N}{\mathbb{N}}
\newcommand{\Z}{\mathbb{Z}}
\newcommand{\Q}{\mathbb{Q}}
\newcommand{\R}{\mathbb{R}}
\newcommand{\C}{\mathbb{C}}
\newcommand{\F}{\mathbb{F}}
\newcommand{\A}{\mathbb{A}}
\newcommand{\PP}{\mathbb{P}}
\newcommand{\E}{\mathbb{E}}
\newcommand{\eps}{\varepsilon}
\newcommand{\dd}{\,\mathrm{d}}
\begin{document}
\section*{UnsolvedMath ID: 30006026}
\textbf{Source problem code:} catalogue entry
\subsection{Definitions and mathematical statement}
\textbf{Scope and unresolved hypotheses.} The supplied record combines Question 6 about volume identities with Conjecture 7 about Poisson limits. The identities and volumes referred to in Question 6 are absent. Its request is therefore retained only as: find a unified proof of the referenced identities and determine further volumes satisfying them. A precise statement of this part requires the missing identities and their hypotheses.

For Conjecture 7, let \(g\) tend to infinity through positive integers, and let \(L_g>0\) be a sequence of boundary lengths. Write \(X_{g,L_g}\) for a random hyperbolic surface of genus \(g\) with one boundary of length \(L_g\). The probability law of \(X_{g,L_g}\) is not specified in the supplied extract; the assertions below require that law to be supplied. Let
\[
N_g^{L_g}(a,b)=\#\{\gamma:\gamma\text{ is a simple closed geodesic on }X_{g,L_g},\ a\leq\ell(\gamma)\leq b\}.
\]
The source does not specify whether the boundary geodesic is counted or how geodesics are identified when counting; these conventions remain unresolved.

For fixed \(0<x<y<\infty\), define the candidate Poisson parameter
\[
\lambda(x,y)=\int_x^y\frac{\cosh t-1}{t}\,\dd t.
\]
The following is a reconstruction of the damaged displayed formulas, rather than a fully specified conjecture. The intended range of endpoints is not supplied; positive finite endpoints are used here to display the formulas unambiguously. The proposed assertion is that there exists a continuous increasing function
\[
r:[0,\infty)\longrightarrow(0,\infty),\qquad r(0)=1,\qquad \lim_{u\to\infty}r(u)=\infty,
\]
independent of the boundary-length sequence and interval, such that the following convergence statements hold for every admissible sequence \((L_g)\):
\[
\begin{aligned}
N_g^{L_g}(x,y)&\xrightarrow{\mathrm d}\operatorname{Poisson}(\lambda(x,y))
&&\text{if }\frac{L_g}{g}\longrightarrow0,\\
N_g^{L_g}\bigl(r(c)x,r(c)y\bigr)&\xrightarrow{\mathrm d}\operatorname{Poisson}(\lambda(x,y))
&&\text{if }\frac{L_g}{g}\longrightarrow c\in(0,\infty),\\
N_g^{L_g}\left(\frac{L_g}{12g}x,\frac{L_g}{12g}y\right)&\xrightarrow{\mathrm d}\operatorname{Poisson}(\lambda(x,y))
&&\text{if }\frac{L_g}{g}\longrightarrow\infty.
\end{aligned}
\]
Here \(\xrightarrow{\mathrm d}\) denotes convergence in distribution, and \(\operatorname{Poisson}(\lambda)\) denotes the distribution with probabilities \(e^{-\lambda}\lambda^k/k!\), \(k=0,1,\ldots\). The domain of \(r\), the positive finite range of \(c\), and the universal quantifiers above express the apparent interpolation target and require confirmation from the complete original formulation. The source does not distinguish strictly increasing from nondecreasing. It gives individual count limits; joint convergence or independence of counts on disjoint intervals is not specified. Although the integrand has continuous extension \(0\) at \(t=0\), inclusion of \(x=0\) in the conjecture is unconfirmed.
\subsection{Short English statement}
The proposed conjecture asks for Poisson limits for counts of simple closed geodesics as the genus grows, with length scales depending on whether the boundary length is smaller than, comparable to, or larger than the genus. A single continuous increasing rescaling function should govern the intermediate regime. The sampling law, counting conventions, endpoint range, and referenced volume identities are missing, so this is a provisional statement requiring mathematical review.
\subsection{Sources}
\begin{itemize}
\item Local UnsolvedMath registry, permanent ID \(30006026\), exact source problem code \texttt{catalogue entry}, title: \emph{Poisson Laws across Boundary Regimes of Random Hyperbolic Surfaces}. The supplied record attributes the material to Question 6 and Conjecture 7 in an Oberwolfach report; no individual proposer is supplied. Its proposed year is 2024.
\item Supplied report citation: \emph{Riemann Surfaces: Random, Flat, and Hyperbolic Geometry} (2025), Oberwolfach Reports, pp. 2369--2424, DOI \texttt{10.4171/owr/2024/41}. \url{https://doi.org/10.4171/owr/2024/41}. The extracted passage includes the page marker 2406 and the designation Oberwolfach Report 41/2024. The differing year designations are retained as supplied.
\item Supplied literature source labeled ``Mirzakhani--Petri'': \url{https://arxiv.org/abs/1710.09727}. No further bibliographic details or precise supporting theorem are supplied.
\item Literature assessment dated exactly \texttt{2026-08-22}: the supplied curation labels the problem open and describes the full three-regime interpolation as conjectural, stating that known Poisson laws for closed or restricted regimes do not establish it. This is supplied dated curation, not an independently verified literature assessment.
\item The record labels its statement \texttt{corrected\_verified} and reports that the official report was checked and its notation repaired. These are supplied provenance claims; no independent source verification is asserted here. The background also describes metric maps as hyperbolic surfaces with very large boundary, citing an unidentified reference [5], and refers to Theorems 1 and 2 whose statements are absent.
\end{itemize}
\end{document}
